\chapter*{Essential terms}
\addcontentsline{toc}{chapter}{Essential terms}
\begin{description}
  \item[Control surface] Where the owner interacts: for example, a phone or a laptop client. It
        need not be the machine executing the work.
  \item[Runtime and writing host] The machine running the agent and hosting the local checkout it
        owns. A separate execution host may run an analysis job.
  \item[Session] A bounded task with an owner, current state, and intended result. A client
        conversation may carry it, but durable records preserve it.
  \item[Worktree] A separate Git checkout used to isolate a writer's files. Each writing boundary
        has an explicit owner.
  \item[Contract] The agreed question, scope, deliverables, validation, and applicable review and
        handoff conditions.
  \item[Oracle] An independent expected result or reference used to judge an implementation. The
        reference's authority and limitations must be known.
  \item[Pinpoint] The exact page, section, equation, figure, table, or record supporting a
        citation at its point of use.
  \item[Handback] Evidence returned by a worker or monitor. The parent still owns any remaining
        analysis, review, and delivery obligations.
  \item[Review gate] A named check required before a claim or artifact can be accepted. Copy,
        source, domain, and rendered-page checks answer different questions.
  \item[Human review] The owner's judgment of the delivered artifact or scientific decision. A
        report's PDF acceptance is distinct from its technical checks.
\end{description}

\section*{Reading the diagrams and labels}

Diagrams use four colours, each with one meaning; the text inside a box always names the
step, so colour is a second cue, never the only one.

\begin{center}
  \begin{tikzpicture}[node distance=4mm]
    \node[work] (w) {Work an agent does};
    \node[check,right=of w] (c) {A check or review};
    \node[result,right=of c] (r) {A result};
    \node[open,right=of r] (o) {A step for a person};
  \end{tikzpicture}
\end{center}

Boxes in the text have fixed roles:
\begin{itemize}
  \item a grey \textbf{headline} bar opens a section with its point;
  \item an orange \textbf{choice} box lists the options for a decision, with a recommendation;
  \item a red \textbf{rule} box states a limit that holds without exception.
\end{itemize}
In the worked examples, \contract, \evidence, \gap\ and \review\ mark requirements,
observations, missing evidence and checks.
